% ============================================================================
%  Research note: the kinetic-relation stability criterion for traveling kinks
%  in damped, driven lattices -- a Jordan-chain identity.
%
%  Standalone; compiles with pdflatex on Overleaf (TeX Live 2023+).
%  Figures are generated by the companion repository into ./figures/ .
%  If a figure file is absent, a labelled placeholder box is drawn instead,
%  so the document always compiles.
%
%  >>> FILL IN the two macros below before circulating. <<<
% ============================================================================
\newcommand{\NoteAuthor}{[Author name]}
\newcommand{\NoteAffil}{[Affiliation]}
\newcommand{\NoteEmail}{[email]}
\newcommand{\RepoURL}{\url{https://github.com/[user]/damped-lattice-kinetic}}
\newcommand{\VKarXiv}{arXiv:[identifier]}

\documentclass[11pt]{article}
\usepackage[T1]{fontenc}
\usepackage{lmodern}
\usepackage{microtype}
\usepackage[margin=1in]{geometry}
\usepackage{amsmath,amssymb,amsthm,mathtools}
\usepackage{bm}
\usepackage{graphicx}
\usepackage{booktabs}
\usepackage{array}
\usepackage{xcolor}
\usepackage{enumitem}
\usepackage[font=small,labelfont=bf]{caption}
\usepackage[colorlinks=true,linkcolor=blue!50!black,citecolor=blue!50!black,urlcolor=blue!50!black]{hyperref}

\graphicspath{{figures/}}

% ---------------------------------------------------------------- figures ---
\newcommand{\figinclude}[2][width=\linewidth]{%
  \IfFileExists{figures/#2}{\includegraphics[#1]{#2}}{%
    \fbox{\parbox[c][0.22\textheight][c]{0.92\linewidth}{\centering\small
      Figure file \texttt{figures/\detokenize{#2}} not yet generated.\\[2pt]
      Run \texttt{make figures} in the companion repository.}}}}

% --------------------------------------------------------------- theorems ---
\theoremstyle{plain}
\newtheorem{theorem}{Theorem}[section]
\newtheorem{proposition}[theorem]{Proposition}
\newtheorem{lemma}[theorem]{Lemma}
\newtheorem{corollary}[theorem]{Corollary}
\theoremstyle{definition}
\newtheorem{hypothesis}{Hypothesis}
\renewcommand{\thehypothesis}{H\arabic{hypothesis}}
\newtheorem{observation}[theorem]{Numerical observation}
\newtheorem{openproblem}{Open problem}
\theoremstyle{remark}
\newtheorem{remark}[theorem]{Remark}

% ----------------------------------------------------------------- macros ---
\newcommand{\ip}[2]{\langle #1,\,#2\rangle}
\newcommand{\ph}{\hat p}
\newcommand{\dd}{\mathrm{d}}
\newcommand{\R}{\mathbb{R}}
\newcommand{\Z}{\mathbb{Z}}
\newcommand{\one}{\mathbf 1}
\newcommand{\T}{\mathsf T}
\DeclareMathOperator{\sech}{sech}
\newcommand{\chat}{\hat c_1}
\newcommand{\cmax}{c_{\max}}
\newcommand{\tagP}{\textsf{proved}}
\newcommand{\tagC}[1]{\textsf{proved under #1}}
\newcommand{\tagN}{\textsf{numerical}}
\newcommand{\tagO}{\textsf{open}}

\title{\textbf{The kinetic-relation stability criterion for traveling kinks\\
in damped, driven lattices: a Jordan-chain identity}\\[6pt]
\large Research note --- draft for discussion}
\author{\NoteAuthor\\ \small \NoteAffil\\ \small\texttt{\NoteEmail}}
\date{Draft of \today}

\begin{document}
\maketitle

\begin{abstract}
In a damped, dc-driven Frenkel--Kontorova lattice, Vainchtein, Cuevas-Maraver,
Kevrekidis and Xu found numerically that traveling kinks acquire an unstable
direction exactly where the kinetic relation $\sigma(c)$ between driving force and
velocity has an extremum, and identified the obstruction to a proof: dissipation
destroys the generalized eigenvector at the neutral Floquet multiplier on which the
Hamiltonian argument rests~\cite{VCKX20}. We show that this obstruction can be
measured exactly. Writing the co-traveling linearization as a quadratic pencil
$Q(\nu)=\nu^2+\nu M_1+M_0$, we prove the operator duality $Q(-\gamma-\nu)=Q(\nu)^{\T}$,
which contains the conformally symplectic structure of the monodromy and identifies the
left neutral eigenvector with the null vector $\ph$ of the \emph{anti-damped}
advance--delay equation. Under hypotheses on $\ph$ that we verify numerically but do not
prove, the scalar $\kappa(c)=\ip{\ph}{M_1\phi'}$ that decides whether the neutral
eigenvalue is simple satisfies $\kappa(c)=-\mu\,\sigma'(c)\,\ip{\ph}{\one}$, and, along the
branch parametrized by arclength, $\dot c\,\kappa+\mu\,\dot\sigma\,\ip{\ph}{\one}=0$.
Hence a multiplier can reach $+1$ only at an extremum of $\sigma$, while at a turning point
of the branch in $c$ it is $\ip{\ph}{\one}$ that vanishes and no crossing occurs. A
Lyapunov--Schmidt reduction then reduces the local sufficiency of the slope criterion to
the sign of one computable coefficient. A Fourier-spectral solver for the advance--delay
equation, independent of the method of~\cite{VCKX20}, reproduces its critical velocity,
critical force and maximal velocity, verifies the identity to $10^{-14}$ with spectral
convergence, and confirms the predicted behaviour at the first extremum and the first
turning point. The identity extends to damped, additively driven lattices with arbitrary
periodic onsite potential and symmetric finite-range coupling. We state precisely what is
proved, what is numerical, and what remains open.
\end{abstract}

\tableofcontents

% ============================================================================
\section{Introduction}\label{sec:intro}
% ============================================================================

\subsection{The problem}

We consider the damped, dc-driven Frenkel--Kontorova (FK) lattice
\begin{equation}\label{eq:FK}
  \ddot u_n+\gamma\dot u_n=u_{n+1}-2u_n+u_{n-1}+\mu\bigl(\sigma-\sin u_n\bigr),
  \qquad n\in\Z,
\end{equation}
with damping $\gamma>0$, onsite strength $\mu>0$ and driving force $\sigma\in(0,1)$, and
its traveling kinks $u_n(t)=\phi(\xi)$, $\xi=n-ct$, which satisfy the advance--delay
equation
\begin{equation}\label{eq:TW}
  c^2\phi''(\xi)-\gamma c\,\phi'(\xi)
  =\phi(\xi+1)-2\phi(\xi)+\phi(\xi-1)+\mu\bigl(\sigma-\sin\phi(\xi)\bigr),
\end{equation}
\begin{equation}\label{eq:BC}
  \phi(\xi)\to\arcsin\sigma+2\pi\ \ (\xi\to-\infty),\qquad
  \phi(\xi)\to\arcsin\sigma\ \ (\xi\to+\infty),\qquad \phi(0)=\pi .
\end{equation}
For $\gamma>0$ the traveling waves come in one-parameter families along which the
driving force is a function of the velocity, the \emph{kinetic relation} $\sigma=\sigma(c)$.
This relation is the central constitutive object in the continuum description of lattice
defects (dislocations, fronts, domain walls) and has been computed repeatedly, together
with the observation that its decreasing portions are unstable
\cite{EW77,WvZSO96,SE98,BHZ00,Car04}.

Vainchtein, Cuevas-Maraver, Kevrekidis and Xu~\cite{VCKX20} computed the traveling kinks
of~\eqref{eq:FK} by Newton iteration on the period map of the lattice, determined their
Floquet multipliers, and found a precise numerical criterion: portions of the kinetic curve
with $\sigma'(c)<0$ always carry an unstable direction, a real multiplier leaving the unit
circle through $+1$ exactly at the extrema of $\sigma$, and a further unstable direction
appears at every subsequent extremum of the spiralling, multivalued kinetic curve. They
note that the sharpness of the criterion makes it ``highly likely to be associated with a
theorem'', and report that the Hamiltonian argument~\cite{FP04,CKVX17,XCKV18} fails for two
reasons. The first is technical: the waves decay slowly to a nontrivial asymptotic state, so
the function spaces and inner products must be adapted. The second is conceptual: in the
Hamiltonian problem the multiplier $\rho=1$ carries both the translation eigenvector
$\partial_\xi\phi$ and the generalized eigenvector $\partial_c\phi$, and the criterion
follows from symplectic orthogonality against the latter; with damping this generalized
eigendirection is eliminated, its multiplier moving to $e^{-\gamma/c}$. They suggest that a
suitably weighted space might restore the Hamiltonian frame.

\subsection{Summary of results}

The present note quantifies the second obstruction exactly and locates the first. Its
results, with their status, are as follows.

\begin{enumerate}[leftmargin=2.2em,label=(\roman*)]
\item \emph{Structure of the linearization} (\S\ref{sec:structure}, \tagP). The monodromy
  $M$ of the linearized lattice is conformally symplectic, $M^{\T}JM=e^{-\gamma/c}J$. The
  co-traveling pencil satisfies the stronger operator identity $Q(-\gamma-\nu)=Q(\nu)^{\T}$,
  which gives the multiplier pairing $\rho\hat\rho=e^{-\gamma/c}$, places the essential
  spectrum on $\operatorname{Re}\nu=-\gamma/2$, and identifies the left neutral eigenvector
  with the null vector $\ph$ of the anti-damped equation ($\gamma\to-\gamma$). The weighted
  space suggested in~\cite{VCKX20} exists in the form of the scalar rescaling
  $e^{\gamma/(2c)}M$, which is exactly symplectic, but it moves the stability boundary off the
  unit circle and does not by itself restore the Hamiltonian argument
  (Remark~\ref{rem:weights}).
\item \emph{The Jordan-chain identity} (\S\ref{sec:identity},
  \tagC{Hypotheses~\ref{H:branch}--\ref{H:adjoint}}). The scalar
  $\kappa(c)=\ip{\ph}{M_1\phi'}$, $M_1=-2c\,\partial_\xi+\gamma$, whose vanishing is
  equivalent to the existence of a generalized eigenvector at $\nu=0$, satisfies
  \begin{equation}\label{eq:main}
    \kappa(c)=-\mu\,\sigma'(c)\,\ip{\ph}{\one},\qquad \ip{\ph}{\cos\phi}=0 .
  \end{equation}
  With the normalization $\ip{\ph}{\one}=\ip{\phi'}{\one}=-2\pi$ this reads
  $\kappa=2\pi\mu\,\sigma'(c)$: the generalized eigendirection is destroyed by the damping by
  exactly $2\pi\mu\sigma'(c)$, and is restored, as $-\partial_c\phi$, precisely where
  $\sigma'(c)=0$.
\item \emph{Turning points} (\S\ref{sec:folds},
  \tagC{Hypotheses~\ref{H:branch}--\ref{H:adjoint}}). Along the branch parametrized by
  arclength, $\dot c\,\kappa+\mu\,\dot\sigma\,\ip{\ph}{\one}=0$. At an extremum of $\sigma$
  this forces $\kappa=0$; at a turning point in $c$ it forces $\ip{\ph}{\one}=0$ instead. The
  two kinds of critical point on the multivalued kinetic curve are thus distinguished by
  which factor of~\eqref{eq:main} vanishes.
\item \emph{Local crossing} (\S\ref{sec:crossing},
  \tagC{Hypotheses~\ref{H:branch}--\ref{H:adjoint} and three non-degeneracy conditions}).
  Near an extremum $\hat c$, the pencil has exactly two eigenvalues near $0$: $\nu=0$ and a
  real $\nu_2(c)$ with $\nu_2\approx-\kappa/m$ and
  $\nu_2'(\hat c)=\mu\sigma''(\hat c)\ip{\ph}{\one}/m(\hat c)$, where $m$ is an explicit
  coefficient whose vanishing at $\hat c$ is equivalent to a Jordan chain of length three.
  The slope criterion holds near $\hat c$ if and only if $m(\hat c)$ has the sign fixed in
  Corollary~\ref{cor:local}.
\item \emph{Power balance} (\S\ref{sec:identity}, \tagP{} for decaying waves).
  $\mu\sigma=(\gamma c/2\pi)\int_\R\phi'(\xi)^2\,\dd\xi$.
\item \emph{Generality} (\S\ref{sec:general}, \tagC{the analogous hypotheses}).
  The identity holds for damped lattices with additive drive $f$, any smooth $2\pi$-periodic
  onsite potential and any symmetric finite-range coupling: $\kappa=-f'(c)\ip{\ph}{\one}$.
\item \emph{Numerics} (\S\ref{sec:numerics}--\ref{sec:results}, \tagN). A Fourier-spectral
  solver for~\eqref{eq:TW} reproduces $\chat$, $\sigma(\chat)$ and $\cmax$ of~\cite{VCKX20};
  confirms~\eqref{eq:main} to $10^{-14}$ with spectral convergence and independence of the
  domain size; confirms that the stability transition and the zero of $\sigma'$ coincide to
  $10^{-6}$ in $c$; confirms the turning-point dichotomy at the first fold; finds $\ph$
  exponentially localized; and confirms the generalized identity in four modified lattices.
\end{enumerate}

What is \emph{not} established is listed in \S\ref{sec:open}. The main gaps are the
functional-analytic Hypotheses~\ref{H:fredholm}--\ref{H:adjoint} (simplicity of the kernel and
exponential decay of $\ph$ on the infinite lattice; the Fredholm property itself is proved in
Lemma~\ref{lem:fredholm}), the sign of $m$, any statement beyond a neighbourhood of each
extremum, crossings through $\rho=-1$ or off the real axis, and the weakly damped resonance
regime.

\subsection{Relation to other stability criteria}

For Hamiltonian lattices the analogous criterion is the change of monotonicity of the
energy along the family, for traveling waves~\cite{FP04,CKVX17,XCKV18} and for discrete
breathers~\cite{KCP16}. In those settings $\partial_c\phi$ is a generalized eigenvector for
every $c$ and the criterion lives one Jordan level higher. The present structure is also
dual to that of the Vakhitov--Kolokolov criterion for solitons of non-Hermitian nonlinear
Schr\"odinger equations~\cite{VK-NH}: there the $U(1)$ Jordan block survives any gain--loss
profile and the relevant scalar is one level up, whereas here the damping destroys
translation's partner and the scalar sits at the first level.

\subsection{Status conventions}
Every statement below carries one of the tags \tagP, \tagC{\dots}, \tagN{} or \tagO. A
\tagN{} statement is supported by convergence-checked computation but not proved; a
statement proved under hypotheses is not claimed beyond them.

% ============================================================================
\section{Setting}\label{sec:setting}
% ============================================================================

\subsection{Linearization, monodromy and the co-traveling pencil}

Linearizing~\eqref{eq:FK} about $u_n(t)=\phi(n-ct)$ gives
\begin{equation}\label{eq:lin}
  \ddot\xi_n+\gamma\dot\xi_n=\xi_{n+1}-2\xi_n+\xi_{n-1}-\mu\cos\bigl(u_n(t)\bigr)\,\xi_n .
\end{equation}
With $z=(\xi,\eta)$, $\eta=\dot\xi$, this is $\dot z=A(t)z$,
\begin{equation}\label{eq:A}
  A(t)=\begin{pmatrix}0&I\\K(t)&-\gamma I\end{pmatrix},\qquad
  K(t)=\Delta-\mu\,\operatorname{diag}\bigl(\cos u_n(t)\bigr)=K(t)^{\T},
\end{equation}
with $\Delta$ the lattice Laplacian. Following~\cite{CKVX17,VCKX20}, the traveling wave is a
fixed point of ``evolve for $T=1/c$, then shift by one site'', and its spectral stability is
governed by the monodromy
\begin{equation}\label{eq:monodromy}
  M=S\,\Phi(T),\qquad (Sz)_n=z_{n+1},
\end{equation}
where $\Phi$ is the fundamental matrix of $\dot z=A(t)z$. On a finite lattice with the
charge-preserving boundary conditions of~\cite{VCKX20}, the perturbation $\xi$ is cyclic and
$S$ is a permutation.

Seeking Floquet solutions $\xi_n(t)=e^{\nu t}p(n-ct)$ of~\eqref{eq:lin} gives the quadratic
eigenvalue problem (Appendix~\ref{app:pencil})
\begin{equation}\label{eq:pencil}
  Q(\nu)p:=\bigl(\nu^2+\nu M_1+M_0\bigr)p=0,\qquad \rho=e^{\nu/c},
\end{equation}
\begin{equation}\label{eq:M0M1}
  M_0=c^2\partial_\xi^2-\gamma c\,\partial_\xi-\Delta_1+\mu\cos\phi,\qquad
  M_1=-2c\,\partial_\xi+\gamma,\qquad
  \Delta_1 p=p(\cdot+1)+p(\cdot-1)-2p .
\end{equation}
$M_0$ is the linearization of~\eqref{eq:TW} at fixed $(\sigma,c)$. The map $\nu\mapsto\rho$ is
$2\pi i c$-periodic; near the real axis this plays no role. Throughout, instability means an
eigenvalue with $\operatorname{Re}\nu>0$, i.e.\ a multiplier with $|\rho|>1$; we rely
on~\cite{CKVX17,XCKV18} for the correspondence between the pencil and the monodromy
formulations.

\subsection{Pairing, transpose, hypotheses}

We use the \emph{bilinear} pairing $\ip fg=\int_\R f(\xi)g(\xi)\,\dd\xi$ and write
$T^{\T}$ for the formal transpose with respect to it. Since
$\partial_\xi^{\T}=-\partial_\xi$ and $\Delta_1^{\T}=\Delta_1$,
\begin{equation}\label{eq:M0T}
  M_0^{\T}=c^2\partial_\xi^2+\gamma c\,\partial_\xi-\Delta_1+\mu\cos\phi,\qquad
  M_1^{\T}=2c\,\partial_\xi+\gamma,
\end{equation}
i.e.\ $M_0^{\T}$ is $M_0$ with $\gamma\to-\gamma$. Let $C^2_b$ denote $C^2$ functions that are
bounded together with their first two derivatives. The following hypotheses are used below.

\begin{hypothesis}[Branch]\label{H:branch}
On an interval $I$ of velocities (or, in \S\ref{sec:folds}, along a $C^1$ curve in
arclength $s$) there is a $C^1$ family $(\phi,\sigma)$ of solutions of~\eqref{eq:TW}--\eqref{eq:BC}
such that $\phi'$, $\phi''$ and $\phi-\phi_{\pm}$ decay exponentially as $\xi\to\pm\infty$
($\phi_\pm$ the asymptotic states), the derivative of the family with respect to the
parameter lies in $C^2_b$, and $\phi'(0)\neq0$.
\end{hypothesis}

\begin{hypothesis}[Simple kernel]\label{H:fredholm}
For each parameter value, $\ker M_0=\operatorname{span}\{\phi'\}$ in $H^2(\R)$.
\end{hypothesis}

\begin{hypothesis}[Adjoint decay]\label{H:adjoint}
The null vector $\ph$ of $M_0^{\T}$ (unique up to scaling under~\ref{H:fredholm}, by
Lemma~\ref{lem:fredholm} below) is $C^2$, and $\ph,\ph',\ph''$ decay exponentially. In
particular $\ph\in L^1$.
\end{hypothesis}

The Fredholm property itself need not be assumed.

\begin{lemma}[Fredholm property; \tagC{\ref{H:branch}}]\label{lem:fredholm}
Let $\gamma>0$ and $\mu\sqrt{1-\sigma^2}>\gamma^2/4$. For every $\nu$ with
$\operatorname{Re}\nu>-\gamma/2$, $Q(\nu):H^2(\R)\to L^2(\R)$ is Fredholm of index zero, and
$\nu\mapsto Q(\nu)$ is analytic there. In particular $M_0=Q(0)$ is Fredholm of index zero; if
\ref{H:fredholm} holds, $\ker M_0^{\T}$ is spanned by a single $\ph\in H^2$ and
$\operatorname{Ran}M_0=\{f\in L^2:\ip{\ph}{f}=0\}$.
\end{lemma}

\begin{proof}
Let $Q_\infty(\nu)$ be $Q(\nu)$ with $\cos\phi$ replaced by its common limit
$\sqrt{1-\sigma^2}$ at $\pm\infty$. It is a Fourier multiplier with symbol
$s(\nu,k)=\nu^2+\nu(\gamma-2ick)+L(k)$, where
$L(k)=-c^2k^2-i\gamma ck+4\sin^2(k/2)+\mu\sqrt{1-\sigma^2}$ (cf.~\cite[eq.~(7)]{VCKX20}).
Writing $\nu=\lambda+ick$ gives $s=\lambda^2+\gamma\lambda+\omega_k^2$ with
$\omega_k^2=4\sin^2(k/2)+\mu\sqrt{1-\sigma^2}>\gamma^2/4$, so $s(\nu,k)=0$ for real $k$ only if
$\operatorname{Re}\nu=-\gamma/2$. For $\operatorname{Re}\nu>-\gamma/2$ the symbol is continuous,
nonvanishing and $\sim-c^2k^2$ as $|k|\to\infty$, hence $|s(\nu,k)|\ge C(1+k^2)$ and
$Q_\infty(\nu):H^2\to L^2$ is an isomorphism. The difference
$Q(\nu)-Q_\infty(\nu)=\mu(\cos\phi-\sqrt{1-\sigma^2})$ is multiplication by a bounded
continuous function tending to zero at $\pm\infty$ by~\ref{H:branch}, which is compact from
$H^2$ to $L^2$. Fredholmness with index zero follows from stability of the index under
compact perturbations. Since the coefficients are real, the $L^2$-adjoint has the same
kernel as the transpose; bootstrapping in $M_0^{\T}\ph=0$ gives $\ph\in H^2$.
\end{proof}

\begin{remark}[Status of the hypotheses]\label{rem:hyp}
Lemma~\ref{lem:fredholm} settles the Fredholm question. What remains is~\ref{H:fredholm},
simplicity of the kernel, which is not automatic, and~\ref{H:adjoint}, exponential decay of
$\ph$ beyond the $L^2$ membership supplied by Lemma~\ref{lem:fredholm}. Since the limiting
operator is hyperbolic ($L(k)\neq0$ for real $k$), exponential decay of $H^2$ kernel elements
is expected from the theory of functional differential equations of mixed
type~\cite{MP99}; we have not carried this out (\tagO). \ref{H:branch} is not addressed
analytically. All three hypotheses are supported numerically (\S\ref{sec:results}).
\end{remark}

% ============================================================================
\section{Structure of the linearized problem}\label{sec:structure}
% ============================================================================

\begin{proposition}[Conformal symplecticity; \tagP]\label{prop:cs}
Let $J=\begin{psmallmatrix}0&I\\-I&0\end{psmallmatrix}$. For any symmetric $K(t)$ in~\eqref{eq:A},
$A^{\T}J+JA=-\gamma J$. Consequently $\Phi(t)^{\T}J\Phi(t)=e^{-\gamma t}J$ and, since $S$ is
symplectic,
\begin{equation}\label{eq:cs}
  M^{\T}JM=e^{-\gamma/c}J .
\end{equation}
\end{proposition}

\begin{proof}
Direct multiplication gives $A^{\T}J=\begin{psmallmatrix}-K&0\\ \gamma I&I\end{psmallmatrix}$ and
$JA=\begin{psmallmatrix}K&-\gamma I\\0&-I\end{psmallmatrix}$, whose sum is
$\begin{psmallmatrix}0&-\gamma I\\ \gamma I&0\end{psmallmatrix}=-\gamma J$. Hence
$\frac{\dd}{\dd t}(\Phi^{\T}J\Phi)=\Phi^{\T}(A^{\T}J+JA)\Phi=-\gamma\,\Phi^{\T}J\Phi$ with
$\Phi(0)=I$. Finally $S=\operatorname{diag}(s,s)$ with $s^{\T}s=I$, so $S^{\T}JS=J$.
\end{proof}

Only three properties of~\eqref{eq:FK} enter: Newtonian second-order form, a symmetric
force Jacobian, and uniform scalar damping. Conformally symplectic structure of damped
Hamiltonian flows is classical~\cite{McLP01}; what matters here is its consequence for
the neutral direction.

\begin{corollary}[Multiplier pairing and left eigenvectors; \tagP]\label{cor:pairing}
The spectrum of $M$ is invariant under $\rho\mapsto e^{-\gamma/c}/\rho$. If $Mv=\rho v$ then
$M^{\T}(Jv)=(e^{-\gamma/c}/\rho)(Jv)$: the left eigenvector at $e^{-\gamma/c}/\rho$ is $Jv$.
In particular the multiplier $\rho=1$ of translation is paired with $\hat\rho=e^{-\gamma/c}$,
recovering the structure reported in~\cite{VCKX20,WvZSO96}.
\end{corollary}

The pencil carries the same information in a form that exhibits the adjoint problem
directly.

\begin{proposition}[Pencil duality; \tagP]\label{prop:duality}
As formal operators on the line, $Q(-\gamma-\nu)=Q(\nu)^{\T}$ for all $\nu\in\mathbb C$.
\end{proposition}

\begin{proof}
Using~\eqref{eq:M0M1},
\begin{align*}
Q(-\gamma-\nu)
&=(\gamma+\nu)^2+(-\gamma-\nu)(-2c\,\partial_\xi+\gamma)+M_0\\
&=\nu^2+\nu\,(2c\,\partial_\xi+\gamma)
 +\bigl[c^2\partial_\xi^2+\gamma c\,\partial_\xi-\Delta_1+\mu\cos\phi\bigr]
 =\nu^2+\nu M_1^{\T}+M_0^{\T},
\end{align*}
where the $\gamma^2$ terms cancel and $2\gamma c\,\partial_\xi-\gamma c\,\partial_\xi=\gamma c\,\partial_\xi$.
By~\eqref{eq:M0T} the right-hand side is $Q(\nu)^{\T}$.
\end{proof}

\begin{corollary}[\tagP]\label{cor:duality}
(a) The spectrum of the pencil is invariant under $\nu\mapsto-\gamma-\nu$, i.e.\
$\rho\hat\rho=e^{-\gamma/c}$. (b) The left null vector of $Q(\nu)$ is the right null vector of
$Q(-\gamma-\nu)$. In particular the left null vector at $\nu=0$ solves $M_0^{\T}\ph=0$, the
anti-damped version of the linearized traveling-wave equation; equivalently, $\ph$
generates the eigenvalue $\nu=-\gamma$, $\hat\rho=e^{-\gamma/c}$. (c) For the
constant-coefficient pencil obtained by replacing $\cos\phi$ by its limit
$\sqrt{1-\sigma^2}$, every spectral value satisfies $\operatorname{Re}\nu=-\gamma/2$, provided
$\mu\sqrt{1-\sigma^2}>\gamma^2/4$.
\end{corollary}

\begin{proof}
(a) and (b) are immediate from Proposition~\ref{prop:duality}. For (c), plane waves
$\xi_n=e^{\lambda t+iqn}$ of~\eqref{eq:lin} with constant coefficients satisfy
$\lambda^2+\gamma\lambda+\omega_q^2=0$, $\omega_q^2=4\sin^2(q/2)+\mu\sqrt{1-\sigma^2}$, so
$\operatorname{Re}\lambda=-\gamma/2$ whenever $\omega_q^2>\gamma^2/4$, which holds for all $q$
under the stated condition; and $\nu=\lambda+icq$ for $p(\xi)=e^{iq\xi}$.
\end{proof}

By Lemma~\ref{lem:fredholm}, the pencil is Fredholm of index zero to the right of
this line, so $|\rho|=e^{-\gamma/(2c)}$ bounds the ``multiplier circle'' of~\cite{VCKX20}
from outside; that the whole line belongs to the essential spectrum follows from standard
Weyl-sequence arguments, which we do not repeat. For the parameters used below
($\mu=1$, $\gamma=0.1$, $\sigma\approx0.65$), $\mu\sqrt{1-\sigma^2}\approx0.76\gg\gamma^2/4$.

\begin{remark}[Two candidate weights; \tagP]\label{rem:weights}
(1) By~\eqref{eq:cs}, $\tilde M=e^{\gamma/(2c)}M$ satisfies $\tilde M^{\T}J\tilde M=J$: a weight
restoring symplecticity exists and is a scalar. It maps the stability boundary $|\rho|=1$ to
$|\tilde\rho|=e^{\gamma/(2c)}\neq1$, and the translation pair $(1,e^{-\gamma/c})$ to the
hyperbolic pair $(e^{\gamma/(2c)},e^{-\gamma/(2c)})$. The Jordan block at $\rho=1$ used in the
Hamiltonian proof is therefore unfolded by an amount depending only on $\gamma/c$, and the
rescaling does not by itself reduce the stability question to a Hamiltonian one.
(2) The substitution $\chi=e^{-\gamma\xi/c}\psi$ intertwines the \emph{differential} parts of
$M_0$ and $M_0^{\T}$ exactly, but maps the shifts to
$e^{\mp\gamma/c}\psi(\xi\pm1)$, so it does not intertwine $M_0$ with $M_0^{\T}$. Numerically,
$\ph$ shows no relative exponential envelope with respect to $\phi'$
(Observation~\ref{obs:decay}). The route taken below bypasses weights altogether and uses
Proposition~\ref{prop:duality}.
\end{remark}

% ============================================================================
\section{The Jordan-chain identity}\label{sec:identity}
% ============================================================================

\begin{lemma}[Pairing lemma; \tagC{\ref{H:adjoint}}]\label{lem:pairing}
If $M_0^{\T}\ph=0$ and $\ph$ satisfies~\ref{H:adjoint}, then $\ip{\ph}{M_0v}=0$ for every
$v\in C^2_b$.
\end{lemma}

\begin{proof}
Integrating by parts, $\int\ph\,c^2v''=\int c^2\ph''v+c^2[\ph v'-\ph'v]_{-\infty}^{\infty}$ and
$\int\ph(-\gamma c\,v')=\int\gamma c\,\ph'v-\gamma c[\ph v]_{-\infty}^{\infty}$; the boundary terms
vanish because $\ph,\ph'\to0$ and $v,v'$ are bounded. The shifts satisfy
$\int\ph(\xi)v(\xi\pm1)\,\dd\xi=\int\ph(\xi\mp1)v(\xi)\,\dd\xi$, all integrals converging
absolutely since $\ph,\ph',\ph''\in L^1$ and $v,v',v''\in L^\infty$. Hence
$\ip{\ph}{M_0v}=\ip{M_0^{\T}\ph}{v}=0$.
\end{proof}

The point of Lemma~\ref{lem:pairing} is that the correct pairing is $L^1$ against
$L^\infty$ rather than $L^2$ against $L^2$: this is what allows pairing against
$\partial_c\phi$, which does not decay, and against the constant function. It is the concrete
form of the ``adapted spaces'' called for in~\cite{VCKX20}.

\begin{lemma}[\tagC{\ref{H:adjoint}}]\label{lem:cos}
$\ip{\ph}{\cos\phi}=0$.
\end{lemma}

\begin{proof}
$M_0\one=\mu\cos\phi$ (derivatives and $\Delta_1$ annihilate constants); apply
Lemma~\ref{lem:pairing} with $v=\one$.
\end{proof}

\begin{proposition}[Power balance; \tagC{\ref{H:branch}}]\label{prop:pb}
Every solution of~\eqref{eq:TW}--\eqref{eq:BC} with exponentially decaying $\phi'$ satisfies
\begin{equation}\label{eq:pb}
  \mu\,\sigma=\frac{\gamma c}{2\pi}\int_\R\phi'(\xi)^2\,\dd\xi .
\end{equation}
\end{proposition}

\begin{proof}
Multiply~\eqref{eq:TW} by $\phi'$ and integrate over $\R$. The term $c^2\phi''\phi'$ is an
exact derivative. For the lattice term,
$\int[\phi(\xi+1)+\phi(\xi-1)]\phi'(\xi)\,\dd\xi=\int\frac{\dd}{\dd\xi}[\phi(\xi)\phi(\xi-1)]\,\dd\xi
=\phi_+^2-\phi_-^2=\int2\phi\phi'$, so $\int(\Delta_1\phi)\phi'=0$. Finally
$\int\mu(\sigma-\sin\phi)\phi'=\mu\sigma(\phi_+-\phi_-)+\mu[\cos\phi]_{-\infty}^{\infty}=-2\pi\mu\sigma$,
since $\phi_\pm$ differ by $2\pi$. What remains is $-\gamma c\int\phi'^2=-2\pi\mu\sigma$.
\end{proof}

Equation~\eqref{eq:pb} expresses the balance between the power injected by the drive and
the power dissipated; the kinetic relation is, up to the factor $\gamma c/2\pi$, the
dissipation functional $\int\phi'^2$. It is the lattice counterpart of the power-balance
admissibility condition for non-Hermitian solitons~\cite{VK-NH}.

\begin{theorem}[Jordan-chain identity; \tagC{\ref{H:branch}, \ref{H:adjoint}}]\label{thm:main}
Along a branch $c\mapsto(\phi,\sigma)$ as in~\ref{H:branch}, with $\ph$ as in~\ref{H:adjoint},
define
\begin{equation}\label{eq:kappa}
  \kappa(c):=\ip{\ph}{M_1\phi'}=\ip{\ph}{-2c\,\phi''+\gamma\,\phi'} .
\end{equation}
Then
\begin{equation}\label{eq:identity}
  \kappa(c)=-\mu\,\sigma'(c)\,\ip{\ph}{\one}.
\end{equation}
With the normalization $\ip{\ph}{\one}=\ip{\phi'}{\one}=\phi_+-\phi_-=-2\pi$,
$\kappa(c)=2\pi\mu\,\sigma'(c)$.
\end{theorem}

\begin{proof}
Write~\eqref{eq:TW} as $F(\phi,\sigma,c)=0$ with
$F=c^2\phi''-\gamma c\phi'-\Delta_1\phi-\mu\sigma+\mu\sin\phi$. Differentiating along the
branch, and using $\partial_cF=2c\phi''-\gamma\phi'=-M_1\phi'$,
\begin{equation}\label{eq:dc}
  M_0[\partial_c\phi]=M_1\phi'+\mu\,\sigma'(c)\,\one .
\end{equation}
By~\ref{H:branch}, $\partial_c\phi\in C^2_b$; note that it does not decay, tending to
$\sigma'(c)/\sqrt{1-\sigma^2}$ at both ends. Pairing~\eqref{eq:dc} with $\ph$ and applying
Lemma~\ref{lem:pairing} gives $0=\kappa+\mu\sigma'\ip{\ph}{\one}$.
\end{proof}

\begin{corollary}[Multiplicity of the neutral eigenvalue; \tagC{\ref{H:branch}--\ref{H:adjoint}}]\label{cor:mult}
$\nu=0$ is an eigenvalue of~\eqref{eq:pencil} with eigenfunction $\phi'\in H^2$. It admits a
Jordan chain of length at least two, i.e.\ there is $p_1\in H^2$ with $M_0p_1+M_1\phi'=0$,
if and only if $\kappa(c)=0$, i.e.\ if and only if $\sigma'(c)\ip{\ph}{\one}=0$. If
$\sigma'(c)=0$, then $\partial_c\phi\in H^2$ and $p_1=-\partial_c\phi$.
\end{corollary}

\begin{proof}
$M_0\phi'=0$ follows by differentiating~\eqref{eq:TW} in $\xi$. By~\ref{H:fredholm} and Lemma~\ref{lem:fredholm}, the
equation $M_0p_1=-M_1\phi'$ (with $M_1\phi'\in L^2$) is solvable iff
$\ip{\ph}{M_1\phi'}=0$, by Lemma~\ref{lem:fredholm}. If $\sigma'=0$, the asymptotes of $\partial_c\phi$ vanish, so
$\partial_c\phi\in H^2$ by~\ref{H:branch}, and~\eqref{eq:dc} reads $M_0\partial_c\phi=M_1\phi'$.
\end{proof}

The damping is visible in~\eqref{eq:kappa} only through the term $\gamma\phi'$ in
$M_1\phi'$ and through $\ph\neq\phi'$. Two ingredients spoil the Hamiltonian chain: the
source $\mu\sigma'\one$ in~\eqref{eq:dc} and the constant tails of $\partial_c\phi$; both are
proportional to $\sigma'(c)$ and vanish together.

\begin{remark}[Hamiltonian limit; \tagP{} as a formal statement]\label{rem:ham}
If $\gamma=0$ and $\phi'$ decays, then $M_0^{\T}=M_0$, $\ph\propto\phi'$, and
$\kappa=-2c\int\phi'\phi''=-c[\phi'^2]_{-\infty}^{\infty}=0$ identically: $\partial_c\phi$ is
then a generalized eigenvector for every $c$, as in~\cite{FP04,XCKV18}. Damping therefore
lowers by one the Jordan level at which the criterion lives. Proposition~\ref{prop:pb}
shows moreover that for $\gamma=0$ a kink with decaying tails forces $\sigma=0$; driven
kinks of the undamped lattice must radiate, consistent with the non-decaying quasiperiodic
tails reported in~\cite{VCKX20}. Since the tail decay rate $|\operatorname{Im}k|$ is
proportional to $\gamma c$, the limit $\gamma\to0$ is singular for the framework of this note.
\end{remark}

% ============================================================================
\section{Turning points: the arclength form and a dichotomy}\label{sec:folds}
% ============================================================================

The kinetic curve of~\eqref{eq:FK} is multivalued~\cite{VCKX20}: $c$ is not a good
parameter at the velocities where branches join.

\begin{theorem}[Arclength identity; \tagC{\ref{H:branch}, \ref{H:adjoint}}]\label{thm:arc}
Let $s\mapsto(\phi(\cdot;s),\sigma(s),c(s))$ be a $C^1$ branch as in~\ref{H:branch}, with
$(\dot\sigma,\dot c)\neq(0,0)$, and let $\kappa(s)=\ip{\ph}{M_1\phi'}$ with $M_1$ evaluated at
$c(s)$. Then
\begin{equation}\label{eq:arc}
  \dot c(s)\,\kappa(s)+\mu\,\dot\sigma(s)\,\ip{\ph}{\one}=0 .
\end{equation}
\end{theorem}

\begin{proof}
Differentiating $F(\phi,\sigma,c)=0$ in $s$ gives
$M_0[\partial_s\phi]=\dot c\,M_1\phi'+\mu\dot\sigma\,\one$; pair with $\ph$ and use
Lemma~\ref{lem:pairing}.
\end{proof}

Equation~\eqref{eq:arc} is homogeneous in $\ph$ and invariant under reparametrization.
Where $\dot c\neq0$ it reduces to~\eqref{eq:identity}.

\begin{corollary}[Dichotomy; \tagC{\ref{H:branch}, \ref{H:adjoint}}]\label{cor:dichotomy}
(a) At an extremum of $\sigma$ with $\dot c\neq0$, $\kappa=0$. (b) At a turning point in $c$
with $\dot\sigma\neq0$, $\ip{\ph}{\one}=0$.
\end{corollary}

At a turning point, \eqref{eq:arc} places no constraint on $\kappa$. Numerically $\kappa$
is bounded away from zero there (Observation~\ref{obs:fold}), so by
Corollary~\ref{cor:mult} the neutral eigenvalue stays simple and no multiplier crosses
$+1$. Table~\ref{tab:dichotomy} summarizes the resulting picture, which accounts for the
observation of~\cite{VCKX20} that new unstable directions appear at the extrema of the
spiralling kinetic curve and not where its branches join.

\begin{table}[h]
\centering\small
\caption{The two kinds of critical point on the multivalued kinetic curve, distinguished by
which factor of $\kappa=-\mu\sigma'\ip{\ph}{\one}$ vanishes. The entries
``$\kappa\neq0$'' and ``no crossing'' in the second row are \tagN.}
\label{tab:dichotomy}
\begin{tabular}{@{}llll@{}}
\toprule
critical point & vanishing factor & $\kappa$ & multiplier at $\rho=1$\\
\midrule
extremum of $\sigma$ ($\dot\sigma=0$, $\dot c\neq0$), e.g.\ $\chat$ & $\sigma'$ & $0$ &
collision; stability changes\\
turning point in $c$ ($\dot c=0$, $\dot\sigma\neq0$), e.g.\ $\cmax$ & $\ip{\ph}{\one}$ &
$\neq0$ & none; stability unchanged\\
\bottomrule
\end{tabular}
\end{table}

\begin{remark}[Discrete counterpart; \tagP{} for the discretization]\label{rem:bordered}
Natural continuation in $c$ solves for $(\psi,\sigma)$ with the bordered Jacobian
$\begin{psmallmatrix}M_0&-\mu\one\\ e_{j_0}^{\T}&0\end{psmallmatrix}$, where $e_{j_0}^{\T}$ is
the pinning functional. When $\ker M_0=\operatorname{span}\{\phi'\}$ and
$\ker M_0^{\T}=\operatorname{span}\{\ph\}$, this matrix is singular if and only if
$\ip{\ph}{\one}=0$ or $\phi'(0)=0$~\cite{Kel77}. The failure of $c$-continuation at a turning
point is therefore the same fact as Corollary~\ref{cor:dichotomy}(b).
\end{remark}

% ============================================================================
\section{Local crossing at an extremum}\label{sec:crossing}
% ============================================================================

Let $\Pi$ denote the projection $f\mapsto\ph\,\ip{\ph}{f}/\ip{\ph}{\ph}$ and let $q_1$ solve
$M_0q_1=(I-\Pi)M_1\phi'$ with $\ip{\phi'}{q_1}=0$. Define
\begin{equation}\label{eq:m}
  m(c):=\ip{\ph}{\phi'}-\ip{\ph}{M_1q_1}.
\end{equation}
At a point $\hat c$ with $\kappa(\hat c)=0$ one has $q_1=\partial_c\phi+\alpha\phi'$ and
$m(\hat c)=\ip{\ph}{\phi'-M_1\partial_c\phi}$.

\begin{proposition}[Local crossing; \tagC{\ref{H:branch}--\ref{H:adjoint}}]\label{prop:crossing}
Let $\sigma'(\hat c)=0$ and assume the non-degeneracy conditions
\[
  \text{(N1) } \sigma''(\hat c)\neq0,\qquad
  \text{(N2) } \ip{\ph}{\one}\neq0 \text{ at }\hat c,\qquad
  \text{(N3) } m(\hat c)\neq0 .
\]
Then there are $\delta,\varepsilon>0$ such that for $|c-\hat c|<\delta$ the pencil has exactly
two eigenvalues, counted with algebraic multiplicity, in $|\nu|<\varepsilon$: $\nu=0$ and a real
eigenvalue $\nu_2(c)$, $C^1$ in $c$, with
\begin{equation}\label{eq:nu2}
  \nu_2(c)=-\frac{\kappa(c)}{m(c)}+O\bigl(\kappa(c)^2\bigr),\qquad
  \nu_2'(\hat c)=\frac{\mu\,\sigma''(\hat c)\,\ip{\ph}{\one}}{m(\hat c)} .
\end{equation}
Moreover (N3) fails if and only if the Jordan chain at $\nu=0$, $\hat c$ has length at least
three.
\end{proposition}

The proof, a Lyapunov--Schmidt reduction, is in Appendix~\ref{app:ls}. The coefficient $m$
plays the role of the length-three non-degeneracy scalar in the analogous non-Hermitian
problem~\cite{VK-NH}.

\begin{corollary}[Local slope criterion; \tagC{Proposition~\ref{prop:crossing}}]\label{cor:local}
Orient $\ph$ so that $\ip{\ph}{\one}<0$ near $\hat c$ (e.g.\ $\ip{\ph}{\one}=-2\pi$). Then
$\operatorname{sign}\nu_2(c)=-\operatorname{sign}\sigma'(c)\cdot\operatorname{sign}m(\hat c)$
for $c$ near $\hat c$. Hence ``$\sigma'<0\Rightarrow$ unstable, $\sigma'>0\Rightarrow$
$\nu_2<0$'' holds near $\hat c$ if and only if $m(\hat c)>0$ in this orientation.
\end{corollary}

Corollary~\ref{cor:local} converts the \emph{necessary} statement of
Corollary~\ref{cor:mult} (a real crossing of $\rho=1$ can occur only where $\kappa=0$) into a
local \emph{sufficient} one, conditional on the sign of a single scalar. It says nothing
about eigenvalues away from $\nu=0$: crossings of $|\rho|=1$ through $\rho=-1$ or by complex
pairs are not addressed.

% ============================================================================
\section{General damped, driven lattices}\label{sec:general}
% ============================================================================

Consider
\begin{equation}\label{eq:gen}
  \ddot u_n+\gamma\dot u_n=\sum_{j=1}^{r}C_j\bigl(u_{n+j}-2u_n+u_{n-j}\bigr)+f-V'(u_n),
\end{equation}
with $V$ smooth and $2\pi$-periodic, finite range $r$, and additive drive $f$. The
traveling-wave operator is
$M_0=c^2\partial_\xi^2-\gamma c\,\partial_\xi-\sum_jC_j\Delta_j+V''(\phi)$, with the same
$M_1=-2c\,\partial_\xi+\gamma$ and $\Delta_jp=p(\cdot+j)+p(\cdot-j)-2p$.

\begin{proposition}[\tagC{the analogues of \ref{H:branch}, \ref{H:adjoint}}]\label{prop:gen}
For a $2\pi$-kink of~\eqref{eq:gen}, with $\ph$ spanning $\ker M_0^{\T}$:
$\kappa=-f'(c)\ip{\ph}{\one}$, $\ip{\ph}{V''(\phi)}=0$, and
$f=(\gamma c/2\pi)\int\phi'^2$. Proposition~\ref{prop:cs}, Proposition~\ref{prop:duality} and
Theorem~\ref{thm:arc} (with $\mu\sigma$ replaced by $f$) hold verbatim.
\end{proposition}

\begin{proof}
Identical to the proofs above: $\partial_cF=-M_1\phi'$, $\partial_fF=-\one$,
$M_0\one=V''(\phi)$, each $\Delta_j$ is symmetric, $\int(\Delta_j\phi)\phi'=0$ by the same
telescoping, and $\int V'(\phi)\phi'=V(\phi_+)-V(\phi_-)=0$.
\end{proof}

Uniform damping is essential for Proposition~\ref{prop:cs}; non-uniform or nonlinear
damping is not covered.

% ============================================================================
\section{Numerical method}\label{sec:numerics}
% ============================================================================

All computations solve the advance--delay equation~\eqref{eq:TW} directly, rather than
iterating the lattice period map as in~\cite{VCKX20}; agreement with~\cite{VCKX20} is
therefore an independent check. Parameters are $\mu=1$, $\gamma=0.1$ unless stated.

\paragraph{Discretization.} Fourier spectral collocation~\cite{Tre00} on $\xi\in[-L,L)$ with
$N$ points, $h=2L/N$, $\xi_j=-L+jh$, and $\xi_{N/2}=0$. Both $\partial_\xi$ and the unit shift
are diagonal in Fourier space, with symbols $ik$ (Nyquist mode set to zero) and $e^{\pm ik}$,
so $c^2\partial_\xi^2-\gamma c\partial_\xi-\Delta_1$ and $M_1$ are exact circulants. The
$2\pi$ jump is carried by the template $A(\xi)=2\pi-4\arctan(e^{\xi/w})$ with $w=2$ fixed;
$A'$, $A''$ and $\Delta_1A$ are evaluated in closed form and the unknown is the periodic
remainder $\psi=\phi-A$, which also carries the constant $\arcsin\sigma$. Discrete pairings
are $\ip fg_h=h\sum_jf_jg_j$.

\paragraph{Traveling waves.} Unknowns $(\psi,\sigma)\in\R^{N+1}$; equations: the $N$
collocated residuals of~\eqref{eq:TW} and the pinning condition $\psi(0)=0$
($\Leftrightarrow\phi(0)=\pi$). Damped Newton with the exact Jacobian
$\begin{psmallmatrix}M_0&-\mu\one\\e_{j_0}^{\T}&0\end{psmallmatrix}$, backtracking by halving
until the max-norm residual decreases (minimum step $10^{-3}$), tolerance $10^{-12}$ on
$\max(\|R\|_\infty,|\psi(0)|)$; achieved residuals are $10^{-13}$. Initialization at
$c_0=0.88$: converge from $\psi=0$, $\sigma=0.6$ with the physical-width template
$w_0=\sqrt{(1-c_0^2)/\mu}$, transfer the profile to the $w=2$ template, re-converge.
Continuation in $c$ with steps $10^{-3}$ (refined to $10^{-4}$ near $\chat$); through turning
points, pseudo-arclength continuation~\cite{Kel77} in $(\psi,\sigma,c)$ with step
$\Delta s=0.02$--$0.03$, tangent from the bordered system, corrector tolerance $10^{-10}$.

\paragraph{Derivatives and null vectors.} $\phi'=A'+\partial_\xi\psi$ spectrally.
$\sigma'(c)$ is computed exactly, not by differencing, from the bordered linear system
$M_0\partial_c\psi-\mu\sigma'\one=M_1\phi'$, $\partial_c\psi(0)=0$. The left null vector $\ph$
is computed by 70--80 steps of inverse iteration on $M_0^{\T}+10^{-14}I$ from a seeded
Gaussian start vector, then normalized either by $\ip{\ph}{\one}_h=-2\pi$ or by
$\|\ph\|_2=1$ with $\ip{\ph}{\phi'}_h>0$. $m$ is computed from~\eqref{eq:m} with $q_1$ from
a bordered solve.

\paragraph{Eigenvalues.} Shift-invert Arnoldi~\cite{LSY98} on the companion linearization
$\begin{psmallmatrix}0&I\\-M_0&-M_1\end{psmallmatrix}$ of~\eqref{eq:pencil}; each application
of the shifted inverse requires one solve with $Q(s)$ via
$x_1=-Q(s)^{-1}\bigl(b_2+(M_1+s)b_1\bigr)$, $x_2=b_1+sx_1$. Typically 14--24 eigenvalues
near a real shift $s$ chosen within $0.004$--$0.035$ of the expected $\nu_2$, tolerance
$10^{-10}$. An eigenvalue is classified as real if $|\operatorname{Im}\nu|<10^{-7}$; the exact
eigenvalues $0$ and $-\gamma$ are identified by $|\nu|<10^{-6}$, $|\nu+\gamma|<10^{-6}$.
Multipliers are $\rho=e^{\nu/c}$.

\paragraph{Resolution.} Production runs use $L=200$ and $N=4096$ ($h\approx0.098$);
continuation through the fold and the model survey use $N=2048$ ($h\approx0.195$). Dense
linear algebra makes the cost $O(N^3)$. At $c=0.88$, $L=200$, the computed asymptotes differ
from $\arcsin\sigma$ and $\arcsin\sigma+2\pi$ by $5\times10^{-9}$, consistent with the tail
decay rate $|\operatorname{Im}k|\approx0.09$ from $L(k)=0$.

% ============================================================================
\section{Numerical results}\label{sec:results}
% ============================================================================

\subsection{Validation against the lattice computation of~\cite{VCKX20}}

\begin{observation}\label{obs:valid}
The values in Table~\ref{tab:valid} agree with~\cite{VCKX20} to all digits reported there.
\end{observation}

\begin{table}[h]
\centering\small
\caption{Critical values on the first branch, $\mu=1$, $\gamma=0.1$. $\chat$ is the zero of the
exactly computed $\sigma'(c)$ (secant iteration, $L=200$, $N=4096$); $\cmax$ is the zero of
$\dot c$ on the pseudo-arclength branch ($N=2048$, linear interpolation).}
\label{tab:valid}
\begin{tabular}{@{}lll@{}}
\toprule
quantity & this note & \cite{VCKX20}\\
\midrule
velocity of the maximum of $\sigma$, $\chat$ & $0.8989297$ & $0.8989$\\
maximal force $\sigma(\chat)$ & $0.6501918$ & $0.65019$\\
maximal velocity (turning point), $\cmax$ & $0.900196$ & $0.9002$\\
\bottomrule
\end{tabular}
\end{table}

\begin{figure}[t]
\centering
\figinclude{fig1_kinetic_relation.pdf}
\caption{Kinetic relation $\sigma(c)$ for $\mu=1$, $\gamma=0.1$. Left: first branch on
$c\in[0.2,0.9]$. Right: neighbourhood of $\chat$ and $\cmax$, including the start of the
second branch obtained by pseudo-arclength continuation; the number of real eigenvalues
with $\operatorname{Re}\nu>0$ is indicated.}
\label{fig:kinetic}
\end{figure}

\subsection{The identity and its convergence}

\begin{observation}\label{obs:identity}
At $c=0.89$ ($L=200$, $N=4096$): $\sigma=0.6373007$, $\sigma'=1.9251321$,
$\kappa=12.095961$ with $\ip{\ph}{\one}=-2\pi$, so that
$\kappa-2\pi\mu\sigma'$ vanishes to machine precision; $\ip{\ph}{\cos\phi}=1.1\times10^{-13}$;
the power balance~\eqref{eq:pb} holds to relative accuracy $4\times10^{-14}$; and
$\partial_c\phi$ at the domain edge equals $2.498179$, against
$\sigma'/\sqrt{1-\sigma^2}=2.498175$. The identity~\eqref{eq:identity} holds to
$10^{-13}$ or better on every point computed with $N=4096$ between $c=0.5$ and $c=0.90015$.
Its relative error converges spectrally in $h$ and is independent of $L$
(Table~\ref{tab:conv}).
\end{observation}

\begin{table}[h]
\centering\small
\caption{Relative error of $\kappa=-\mu\sigma'\ip{\ph}{\one}$ at $c=0.89$, with $\sigma'$
from the exact bordered system, and residual of the translation mode of the discretized
operator.}
\label{tab:conv}
\begin{tabular}{@{}rrrcc@{}}
\toprule
$L$ & $N$ & $h$ & relative error of~\eqref{eq:identity} & $\|M_0\phi'\|_\infty/\|\phi'\|_\infty$\\
\midrule
200 & 1024 & 0.391 & $3.1\times10^{-4}$ & $9.8\times10^{-2}$\\
200 & 2048 & 0.195 & $1.3\times10^{-8}$ & $2.9\times10^{-4}$\\
200 & 4096 & 0.098 & $\lesssim10^{-14}$ & $2.5\times10^{-9}$\\
300 & 3072 & 0.195 & $1.3\times10^{-8}$ & $2.9\times10^{-4}$\\
300 & 6144 & 0.098 & $1.2\times10^{-13}$ & $2.5\times10^{-9}$\\
150 & 2048 & 0.146 & $2.7\times10^{-12}$ & $6.8\times10^{-6}$\\
\bottomrule
\end{tabular}
\end{table}

The error is not a statement about~\eqref{eq:identity} itself, which is exact for the
continuous problem: the discrete branch equation reproduces~\eqref{eq:dc} only up to the
spectral error in differentiating the analytic template, which is why the error tracks
$\|M_0\phi'\|/\|\phi'\|$. Replacing the exact $\sigma'$ by a centred difference with step
$10^{-3}$ yields a grid-independent discrepancy of $1.6\times10^{-4}$, the differencing error
near the curvature maximum of $\sigma$.

\begin{figure}[t]
\centering
\figinclude{fig2_jordan_identity.pdf}
\caption{Left: $\kappa(c)$ and $2\pi\mu\sigma'(c)$ on the first branch (normalization
$\ip{\ph}{\one}=-2\pi$), with their relative difference in the lower panel. Right: the
largest real nontrivial eigenvalue $\nu_2$ from shift-invert Arnoldi, and the leading-order
prediction $-\kappa/m$ of Proposition~\ref{prop:crossing}, near $\chat$.}
\label{fig:identity}
\end{figure}

\subsection{The first extremum}

\begin{observation}\label{obs:threshold}
Near $\chat$ (Table~\ref{tab:threshold}): (a) $\nu_2$ changes sign where $\sigma'$ does;
by linear interpolation the zeros are $0.8989293$ ($\sigma'$) and $0.8989296$ ($\nu_2$),
against $\chat=0.8989297$ from the secant iteration, so the stability transition and the
extremum coincide to $10^{-6}$ in $c$. (b) $-\kappa/m$ converges to $\nu_2$ as
$c\to\chat$: relative discrepancy $10\%$ at $c=0.898$, $1.5\%$ at $0.8988$, and agreement to the
digits shown at $0.89892$. (c) $m(\chat)\approx283>0$ in the orientation
$\ip{\ph}{\one}=-2\pi$, so by Corollary~\ref{cor:local} the slope criterion holds locally at
$\chat$ (conditionally on its hypotheses). (d) $m$ changes sign between $c=0.8999$ and
$0.9000$ on the unstable side, where $\nu_2\approx0.03$--$0.04$ is comparable to the distance
$\gamma/2=0.05$ to the essential spectrum. There $-\kappa/m$ no longer approximates
$\nu_2$, while $\nu_2$ itself stays positive: the sign change marks the edge of the domain of
validity of the leading-order reduction, not a change of stability. (e) For
$c\in\{0.20,0.25,\dots,0.80\}\cup\{0.850,0.855,\dots,0.895\}$ ($N=2048$), $m>0$ in the same
orientation.
\end{observation}

\begin{table}[h]
\centering\small
\caption{Near the first extremum; $L=200$, $N=4096$, normalization $\ip{\ph}{\one}=-2\pi$
(so $\kappa=2\pi\mu\sigma'$ to the digits shown). $\nu_2$: largest real eigenvalue other than
$0$ and $-\gamma$ (shift-invert Arnoldi); $\rho_2=e^{\nu_2/c}$. At $c=0.895$ the predicted
value lies $0.004$ from the essential-spectrum line and no isolated real eigenvalue was
resolved. At $c=0.9$ the reduction is outside its range of validity.}
\label{tab:threshold}
\begin{tabular}{@{}lrrrrrl@{}}
\toprule
$c$ & $\sigma'(c)$ & $\kappa$ & $m$ & $-\kappa/m$ & $\nu_2$ & $\rho_2$\\
\midrule
0.89500 & $ 1.5331386$ & $  9.6329939$ & $209.0$ & $-0.046091$ & ---         & ---\\
0.89800 & $ 0.7418062$ & $  4.6609061$ & $236.5$ & $-0.019705$ & $-0.021961$ & 0.97584\\
0.89880 & $ 0.1477746$ & $  0.9284954$ & $274.0$ & $-0.003389$ & $-0.003441$ & 0.99618\\
0.89892 & $ 0.0118952$ & $  0.0747397$ & $283.4$ & $-0.000264$ & $-0.000264$ & 0.99971\\
0.89900 & $-0.0900527$ & $ -0.5658180$ & $290.4$ & $+0.001948$ & $+0.001931$ & 1.00215\\
0.89920 & $-0.3976031$ & $ -2.4982139$ & $310.5$ & $+0.008047$ & $+0.007740$ & 1.00864\\
0.89950 & $-1.0916745$ & $ -6.8591933$ & $336.9$ & $+0.020362$ & $+0.017591$ & 1.01975\\
0.89980 & $-2.4785852$ & $-15.573410$  & $245.0$ & $+0.063567$ & $+0.029743$ & 1.03361\\
0.90000 & $-4.8695485$ & $-30.596276$  & $-101.4$& (invalid)   & $+0.040624$ & 1.04617\\
\bottomrule
\end{tabular}
\end{table}

\subsection{The first turning point}

\begin{observation}\label{obs:fold}
Along the pseudo-arclength branch through $\cmax$ ($N=2048$, $\Delta s=0.03$,
normalization $\|\ph\|_2=1$, $\ip{\ph}{\phi'}>0$; Table~\ref{tab:fold}, Figure~\ref{fig:fold}):
(a) $\dot c$ and $\ip{\ph}{\one}$ change sign in the same continuation step, at fractions
$0.336$ and $0.342$ of that step by linear interpolation. (b) $\kappa$ passes the turning
point smoothly and bounded away from zero ($\kappa\approx-2.54$). (c) Exactly one real
eigenvalue with $\operatorname{Re}\nu>0$ is found before, at and after the turning point,
increasing continuously through $\nu\approx0.065$ ($\rho\approx1.075$); no eigenvalue crosses
$\nu=0$. (d) In a separate run with $\Delta s=0.02$, the invariant~\eqref{eq:arc} was evaluated
at every step from $c=0.89984$ through the turning point and remained between
$2\times10^{-9}$ and $8\times10^{-9}$, the discretization level of the $N=2048$ grid
(cf.\ Table~\ref{tab:conv}).
\end{observation}

\begin{table}[h]
\centering\small
\caption{Pseudo-arclength steps across the turning point $\cmax$; $\nu$ is the unique real
eigenvalue with $\operatorname{Re}\nu>0$. The sign convention $\ip{\ph}{\phi'}>0$ flips a
few steps later (both $\kappa$ and $\ip{\ph}{\one}$ change sign together); the
invariant~\eqref{eq:arc} is unaffected. $\ph$ has unit Euclidean norm as a grid vector.}
\label{tab:fold}
\begin{tabular}{@{}rllrrrl@{}}
\toprule
step & $c$ & $\sigma$ & $\dot c$ & $\kappa$ & $\ip{\ph}{\one}$ & $\nu$ ($\rho$)\\
\midrule
21 & 0.900193 & 0.646593 & $ 9.74\times10^{-5}$ & $-2.58417$ & $-0.044287$ & 0.06252 (1.07192)\\
22 & 0.900195 & 0.646421 & $ 5.54\times10^{-5}$ & $-2.56210$ & $-0.024460$ & 0.06392 (1.07359)\\
23 & 0.900196 & 0.646246 & $ 1.38\times10^{-5}$ & $-2.54027$ & $-0.005925$ & 0.06531 (1.07525)\\
24 & 0.900196 & 0.646067 & $-2.73\times10^{-5}$ & $-2.51881$ & $ 0.011415$ & 0.06670 (1.07691)\\
25 & 0.900194 & 0.645884 & $-6.79\times10^{-5}$ & $-2.49779$ & $ 0.027650$ & 0.06808 (1.07856)\\
\bottomrule
\end{tabular}
\end{table}

\begin{figure}[t]
\centering
\figinclude{fig3_fold_dichotomy.pdf}
\caption{Quantities along arclength through the first turning point: $\dot c$,
$\ip{\ph}{\one}$, $\kappa$ and the unstable eigenvalue $\nu$. The zeros of $\dot c$ and
$\ip{\ph}{\one}$ coincide; $\kappa$ and $\nu$ are smooth and nonzero.}
\label{fig:fold}
\end{figure}

\subsection{The adjoint null vector and the spectrum}

\begin{observation}\label{obs:decay}
At $c=0.89$ ($L=200$, $N=4096$), $\|M_0^{\T}\ph\|/\|\ph\|=7\times10^{-14}$. Both null vectors
are localized ($99.0\%$ of the discrete $L^2$ mass in $|\xi|<20$ at $N=2048$) and decay with envelopes of equal
magnitude (Table~\ref{tab:decay}, Figure~\ref{fig:null}). A relative factor
$e^{\pm\gamma\xi/c}$ between them, as a naive exponential weight would imply, would amount to
$e^{\pm22}$ at $|\xi|=199$; none is present. At $N=2048$ the three smallest singular values of $M_0$ are
$2.7\times10^{-9}$, $9.9\times10^{-2}$ and $1.9\times10^{-1}$, consistent with a
one-dimensional kernel (Hypothesis~\ref{H:fredholm}).
\end{observation}

\begin{table}[h]
\centering\small
\caption{$\max(|f(\xi)|,|f(-\xi)|)/\max|f|$ for $f=\phi'$ and $f=\ph$ at $c=0.89$. The tails
are oscillatory, hence not monotone.}
\label{tab:decay}
\begin{tabular}{@{}lrrrrrrr@{}}
\toprule
$|\xi|$ & 0 & 10 & 20 & 40 & 80 & 120 & 199\\
\midrule
$\phi'$ & $6.7\times10^{-1}$ & $1.2\times10^{-1}$ & $5.3\times10^{-3}$ & $1.2\times10^{-2}$ &
$1.2\times10^{-4}$ & $2.7\times10^{-7}$ & $1.1\times10^{-9}$\\
$\ph$ & $4.8\times10^{-1}$ & $1.5\times10^{-1}$ & $6.7\times10^{-3}$ & $1.2\times10^{-2}$ &
$1.2\times10^{-4}$ & $1.2\times10^{-7}$ & $1.1\times10^{-9}$\\
\bottomrule
\end{tabular}
\end{table}

\begin{figure}[t]
\centering
\figinclude{fig4_null_vectors.pdf}
\caption{$|\phi'|$ and $|\ph|$ at $c=0.89$ on a logarithmic scale.}
\label{fig:null}
\end{figure}

\begin{observation}\label{obs:spectrum}
At $c=0.89$ the computed spectrum near the origin contains $\nu=7\times10^{-9}$ and
$\nu=-0.100000007$ (the exact pair $0,-\gamma$) and no other real eigenvalue; every computed
eigenvalue has a partner with $\nu+\hat\nu=-\gamma$ to $10^{-8}$; the remaining eigenvalues
cluster on $\operatorname{Re}\nu=-\gamma/2$, i.e.\ $|\rho|=e^{-\gamma/(2c)}=0.945369$
(Figure~\ref{fig:spectrum}). For a random symmetric $K(t)$ ($8$ sites, $\gamma=0.37$,
$T=0.9$), direct integration gives $\|M^{\T}JM-e^{-\gamma T}J\|/\|J\|=1.9\times10^{-14}$,
confirming Proposition~\ref{prop:cs} numerically.
\end{observation}

\begin{figure}[t]
\centering
\figinclude{fig5_spectrum.pdf}
\caption{Pencil eigenvalues near the origin in the $\nu$-plane (left) and the corresponding
multipliers (right) at $c=0.89$ (stable) and $c=0.8995$ (one unstable real multiplier). The
dashed line $\operatorname{Re}\nu=-\gamma/2$ and circle $|\rho|=e^{-\gamma/(2c)}$ are the axis
of the symmetry $\nu\mapsto-\gamma-\nu$ and the essential spectrum.}
\label{fig:spectrum}
\end{figure}

\subsection{Other lattices}

\begin{observation}\label{obs:gen}
For the lattices of Table~\ref{tab:gen}, all at $c=0.82$, $L=200$, $N=2048$, $\|\ph\|_2=1$,
with $f'(c)$ from the exact bordered system, $\kappa=-f'(c)\ip{\ph}{\one}$ to the accuracy
shown. The larger error for $\mu_2=-0.5$ is a resolution effect: at $N=4096$ it is
$2.4\times10^{-12}$.
\end{observation}

\begin{table}[h]
\centering\small
\caption{Test of Proposition~\ref{prop:gen}. $V'(u)=\mu\sin u+\mu_2\sin2u$; $C_1=1$
throughout; $\gamma=0.1$ unless stated. $\ph$ is normalized in the Euclidean norm of the
grid vector, so $\kappa$ and $\ip{\ph}{\one}$ individually depend on $N$; their relation does
not. A case with $\mu_2=0.35$ was attempted; natural continuation from $\psi=0$ did not
converge and it is not reported.}
\label{tab:gen}
\begin{tabular}{@{}lrrrrr@{}}
\toprule
lattice & $f$ & $f'(c)$ & $\kappa$ & $-f'\ip{\ph}{\one}$ & rel.\ error\\
\midrule
FK, $\mu=1$ & 0.481960 & 2.174156 & 1.146659 & 1.146659 & $1.8\times10^{-10}$\\
$\mu_2=-0.5$ & 0.540326 & 2.676268 & 1.069733 & 1.069737 & $3.9\times10^{-6}$\\
$C_2=0.25$ & 0.166333 & 0.601746 & 0.480627 & 0.480627 & $1.5\times10^{-14}$\\
$\mu_2=0.3$, $C_2=0.2$ & 0.191512 & 1.009012 & 0.735407 & 0.735407 & $2.1\times10^{-11}$\\
$\mu=0.6$, $\gamma=0.2$ & 0.379908 & 1.487688 & 1.177935 & 1.177935 & $1.5\times10^{-13}$\\
\bottomrule
\end{tabular}
\end{table}

\begin{figure}[t]
\centering
\figinclude{fig6_convergence.pdf}
\caption{Relative error of the identity~\eqref{eq:identity} at $c=0.89$ against $h$, for
$L\in\{150,200,300\}$, together with the translation-mode residual.}
\label{fig:conv}
\end{figure}

% ============================================================================
\section{What is established and what is open}\label{sec:open}
% ============================================================================

\begin{table}[h]
\centering\small
\caption{Status of the main statements.}
\label{tab:status}
\begin{tabular}{@{}p{0.52\linewidth}p{0.18\linewidth}p{0.22\linewidth}@{}}
\toprule
statement & status & where\\
\midrule
$M^{\T}JM=e^{-\gamma/c}J$; $Q(-\gamma-\nu)=Q(\nu)^{\T}$; left neutral vector solves
$M_0^{\T}\ph=0$ & \tagP & Props.~\ref{prop:cs}, \ref{prop:duality}\\
scalar rescaling symplectic but not stability-preserving; exponential weight fails on
shifts & \tagP & Remark~\ref{rem:weights}\\
power balance $\mu\sigma=(\gamma c/2\pi)\int\phi'^2$ & \tagC{\ref{H:branch}} & Prop.~\ref{prop:pb}\\
$\kappa=-\mu\sigma'\ip{\ph}{\one}$; $\ip{\ph}{\cos\phi}=0$ &
\tagC{\ref{H:branch},\ref{H:adjoint}} & Thm.~\ref{thm:main}\\
$\nu=0$ has a Jordan chain iff $\kappa=0$ & \tagC{\ref{H:branch}--\ref{H:adjoint}} &
Cor.~\ref{cor:mult}\\
$\dot c\kappa+\mu\dot\sigma\ip{\ph}{\one}=0$; $\ip{\ph}{\one}=0$ at turning points in $c$ &
\tagC{\ref{H:branch},\ref{H:adjoint}} & Thm.~\ref{thm:arc}, Cor.~\ref{cor:dichotomy}\\
local crossing law~\eqref{eq:nu2}; local slope criterion iff $m(\hat c)>0$ &
\tagC{\ref{H:branch}--\ref{H:adjoint}, (N1)--(N3)} & Prop.~\ref{prop:crossing},
Cor.~\ref{cor:local}\\
generalization to~\eqref{eq:gen} & \tagC{analogues} & Prop.~\ref{prop:gen}\\
$Q(\nu)$ Fredholm of index zero for $\operatorname{Re}\nu>-\gamma/2$ &
\tagC{\ref{H:branch}} & Lemma~\ref{lem:fredholm}\\
\ref{H:branch}--\ref{H:adjoint} at $\gamma=0.1$; identity to $10^{-14}$; $\chat$, $\cmax$;
dichotomy at $\cmax$; $m(\chat)>0$ & \tagN & \S\ref{sec:results}\\
\ref{H:fredholm}--\ref{H:adjoint} on the infinite lattice & \tagO & Remark~\ref{rem:hyp}\\
sign of $m(\hat c)$ in general; global criterion & \tagO & below\\
\bottomrule
\end{tabular}
\end{table}

\begin{openproblem}[Fredholm theory]\label{op:fredholm}
Prove~\ref{H:fredholm}--\ref{H:adjoint} for $\gamma>0$: show that $\ker M_0$ is simple along
the branch, and deduce the exponential decay of $\ph$ (hence $\ph\in L^1$) from the
hyperbolicity of the limiting operator~\cite{MP99}. The Fredholm index is already zero by
Lemma~\ref{lem:fredholm}. This is the rigorous form of the ``adapted spaces'' question
of~\cite{VCKX20}, and it converts every conditional statement above into an unconditional
one on the branches where~\ref{H:branch} holds.
\end{openproblem}

\begin{openproblem}[Sign of $m$]\label{op:m}
Prove that $m(\hat c)>0$ (orientation $\ip{\ph}{\one}<0$) at every extremum of $\sigma$, or
characterize the extrema where it fails. At a failure the slope criterion itself would fail
locally, and the Jordan chain would have length at least three when $m(\hat c)=0$. A
structural reason for the sign is not known.
\end{openproblem}

\begin{openproblem}[Beyond the local window]\label{op:global}
Proposition~\ref{prop:crossing} is local. Numerically $\nu_2$ remains positive on the
decreasing part of the first branch and through $\cmax$, but no argument controls it once
$|\nu_2|$ is comparable to $\gamma/2$. Crossings through $\rho=-1$ and by complex pairs are
not addressed by the present analysis.
\end{openproblem}

\begin{openproblem}[Higher extrema]\label{op:higher}
Theorem~\ref{thm:main} predicts $\kappa=0$ at every extremum $\hat c_n$ of the spiralling
kinetic curve reported in~\cite{VCKX20} at which $\ip{\ph}{\one}\neq0$. Only the first
extremum and the first turning point have been computed here.
\end{openproblem}

\begin{openproblem}[Weak damping and resonances]\label{op:resonance}
For $\gamma=0.01$ the kinetic curve of~\cite{VCKX20} is non-monotone at low velocities
because of resonances with linear waves, near $c\approx0.157$. The tail decay rate is
proportional to $\gamma c$ and tends to zero at resonance, so this is the regime
where~\ref{H:adjoint} is weakest and where the identity should be tested next. It requires
domains roughly an order of magnitude longer than those used here.
\end{openproblem}

\begin{openproblem}[Other dissipative lattices]\label{op:other}
Proposition~\ref{prop:gen} covers uniform linear damping. Non-uniform or nonlinear damping
breaks Proposition~\ref{prop:cs}; viscoelastic chains~\cite{Vai25} provide a natural test
case with a different dissipation mechanism and, there, period-doubling instabilities.
\end{openproblem}

% ============================================================================
\section*{Reproducibility}\addcontentsline{toc}{section}{Reproducibility}
% ============================================================================

All numbers, tables and figures in this note are produced by the companion repository
\RepoURL. Each observation is generated by a numbered script, and every quoted number is
listed with its tolerance in a machine-checked manifest (\texttt{manuscript/claims.yaml}),
verified by \texttt{make check}. Parameters, discretizations, tolerances, seeds and solver
settings are recorded with each output. The repository also contains a direct
time-integration check of the lattice monodromy, which is not used in this draft.

\section*{Acknowledgments}\addcontentsline{toc}{section}{Acknowledgments}
[To be completed.]
% Reminder: add any disclosure of AI-assisted computation or writing required by your
% institution or by the eventual venue.

% ============================================================================
\appendix
\section{Derivation of the pencil}\label{app:pencil}
% ============================================================================

With $u_n(t)=\phi(n-ct)$ and $\xi_n(t)=e^{\nu t}p(n-ct)$,
$\dot\xi_n=e^{\nu t}(\nu p-cp')$ and $\ddot\xi_n=e^{\nu t}(\nu^2p-2\nu cp'+c^2p'')$, evaluated at
$\xi=n-ct$. Substituting into~\eqref{eq:lin} and dividing by $e^{\nu t}$ gives
$\nu^2p+\nu(-2cp'+\gamma p)+\bigl(c^2p''-\gamma cp'-\Delta_1p+\mu\cos\phi\,p\bigr)=0$,
which is~\eqref{eq:pencil}. Since $\xi_{n+1}(T)=e^{\nu T}p(n+1-cT)=e^{\nu T}\xi_n(0)$ with
$T=1/c$, the corresponding multiplier of~\eqref{eq:monodromy} is $\rho=e^{\nu/c}$. Setting
$\nu=0$ recovers $M_0\phi'=0$; setting $\nu=-\gamma$ gives $Q(-\gamma)=M_0^{\T}$, so the
multiplier $e^{-\gamma/c}$ is generated by $\ph$.

% ============================================================================
\section{Proof of Proposition~\ref{prop:crossing}}\label{app:ls}
% ============================================================================

Fix $c$ near $\hat c$. By~\ref{H:fredholm} and Lemma~\ref{lem:fredholm}, write
$H^2=\operatorname{span}\{\phi'\}\oplus X_c$ with $X_c=\{w:\ip{\phi'}{w}=0\}$ and
$L^2=\operatorname{span}\{\ph\}\oplus\operatorname{Ran}M_0$. The equation $Q(\nu)(\phi'+w)=0$,
$w\in X_c$, splits into
\[
  (I-\Pi)Q(\nu)(\phi'+w)=0,\qquad \ip{\ph}{Q(\nu)(\phi'+w)}=0 .
\]
Since $(I-\Pi)M_0:X_c\to\operatorname{Ran}M_0$ is invertible, the first equation has a unique
solution $w=w(\nu,c)$, analytic in $\nu$ and $C^1$ in $c$, with $w(0,c)=0$ because
$Q(0)\phi'=M_0\phi'=0$. Expanding, $w=-\nu q_1+O(\nu^2)$ with $q_1$ as in~\eqref{eq:m}. The
bifurcation function is
\[
  b(\nu,c)=\ip{\ph}{Q(\nu)(\phi'+w)}
  =\nu\,\kappa(c)+\nu^2\bigl(\ip{\ph}{\phi'}-\ip{\ph}{M_1q_1}\bigr)+O(\nu^3)
  =\nu\,h(\nu,c),
\]
using $\ip{\ph}{M_0w}=0$, with $h(0,c)=\kappa(c)$ and $\partial_\nu h(0,c)=m(c)$. The
eigenvalues near $0$ are the zeros of $b$, counted with multiplicity by the Gohberg--Sigal
theory of analytic Fredholm operator functions~\cite{GS71}. The factor $\nu$ gives the
translation eigenvalue. At $(0,\hat c)$, $h=\kappa(\hat c)=0$ by Theorem~\ref{thm:main}, and
$\partial_\nu h=m(\hat c)\neq0$ by (N3), so the implicit function theorem gives a unique
$C^1$ zero $\nu_2(c)$ of $h$ near $0$, with $\nu_2=-\kappa/m+O(\kappa^2)$ and
$\nu_2'(\hat c)=-\kappa'(\hat c)/m(\hat c)$. By Theorem~\ref{thm:main},
$\kappa'(\hat c)=-\mu\sigma''(\hat c)\ip{\ph}{\one}$ at $\hat c$ (the term proportional to
$\sigma'(\hat c)$ vanishes), which is nonzero by (N1)--(N2); this gives~\eqref{eq:nu2}.
Because the pencil has real coefficients, $\bar\nu_2$ is also an eigenvalue; uniqueness
forces $\nu_2\in\R$. For the last statement, at $\hat c$ the Keldysh chain is
$p_0=\phi'$, $p_1=-\partial_c\phi$ (Corollary~\ref{cor:mult}); a third vector exists iff
$M_0p_2=-M_1p_1-p_0=M_1\partial_c\phi-\phi'$ is solvable, i.e.\ iff
$\ip{\ph}{\phi'-M_1\partial_c\phi}=m(\hat c)=0$~\cite{Mar88}.

% ============================================================================
\begin{thebibliography}{99}\small

\bibitem{VCKX20}
A.~Vainchtein, J.~Cuevas-Maraver, P.~G. Kevrekidis, H.~Xu,
Stability of traveling waves in a driven Frenkel--Kontorova model,
\emph{Commun. Nonlinear Sci. Numer. Simul.} \textbf{85} (2020) 105236;
arXiv:1912.05052.

\bibitem{CKVX17}
J.~Cuevas-Maraver, P.~G. Kevrekidis, A.~Vainchtein, H.~Xu,
Unifying perspective: solitary traveling waves as discrete breathers in Hamiltonian
lattices and energy criteria for their stability,
\emph{Phys. Rev. E} \textbf{96} (2017) 032214; arXiv:1701.04882.

\bibitem{XCKV18}
H.~Xu, J.~Cuevas-Maraver, P.~G. Kevrekidis, A.~Vainchtein,
An energy-based stability criterion for solitary travelling waves in Hamiltonian
lattices,
\emph{Phil. Trans. R. Soc. A} \textbf{376} (2018) 20170192; arXiv:1711.03330.

\bibitem{KCP16}
P.~G. Kevrekidis, J.~Cuevas-Maraver, D.~E. Pelinovsky,
Energy criterion for the spectral stability of discrete breathers,
\emph{Phys. Rev. Lett.} \textbf{117} (2016) 094101.

\bibitem{FP04}
G.~Friesecke, R.~L. Pego,
Solitary waves on FPU lattices: III. Howland-type Floquet theory,
\emph{Nonlinearity} \textbf{17} (2004) 207--227.

\bibitem{Car04}
A.~Carpio,
Nonlinear stability of oscillatory wave fronts in chains of coupled oscillators,
\emph{Phys. Rev. E} \textbf{69} (2004) 046601.

\bibitem{EW77}
Y.~Y. Earmme, J.~H. Weiner,
Dislocation dynamics in the modified Frenkel--Kontorova model,
\emph{J. Appl. Phys.} \textbf{48} (1977) 3317--3331.

\bibitem{WvZSO96}
S.~Watanabe, H.~S.~J. van der Zant, S.~H. Strogatz, T.~P. Orlando,
Dynamics of circular arrays of Josephson junctions and the discrete sine-Gordon equation,
\emph{Physica D} \textbf{97} (1996) 429--470.

\bibitem{SE98}
T.~Strunz, F.-J. Elmer,
Driven Frenkel--Kontorova model. I. Uniform sliding states and dynamical domains of
different particle densities,
\emph{Phys. Rev. E} \textbf{58} (1998) 1601--1611.

\bibitem{BHZ00}
O.~M. Braun, B.~Hu, A.~Zeltser,
Driven kink in the Frenkel--Kontorova model,
\emph{Phys. Rev. E} \textbf{62} (2000) 4235--4245.

\bibitem{Vai25}
A.~Vainchtein,
Stability of subsonic kinks in a viscoelastic chain,
\emph{Math. Mech. Solids} (2025), doi:10.1177/10812865251344538.

\bibitem{MP99}
J.~Mallet-Paret,
The Fredholm alternative for functional differential equations of mixed type,
\emph{J. Dynam. Differential Equations} \textbf{11} (1999) 1--47.

\bibitem{McLP01}
R.~McLachlan, M.~Perlmutter,
Conformal Hamiltonian systems,
\emph{J. Geom. Phys.} \textbf{39} (2001) 276--300.

\bibitem{GS71}
I.~C. Gohberg, E.~I. Sigal,
An operator generalization of the logarithmic residue theorem and the theorem of
Rouch\'e,
\emph{Math. USSR-Sb.} \textbf{13} (1971) 603--625.

\bibitem{Mar88}
A.~S. Markus,
\emph{Introduction to the Spectral Theory of Polynomial Operator Pencils},
Transl. Math. Monogr.\ \textbf{71}, Amer. Math. Soc., Providence, 1988.

\bibitem{Kel77}
H.~B. Keller,
Numerical solution of bifurcation and nonlinear eigenvalue problems,
in: P.~H. Rabinowitz (ed.), \emph{Applications of Bifurcation Theory},
Academic Press, New York, 1977, pp.~359--384.

\bibitem{Tre00}
L.~N. Trefethen,
\emph{Spectral Methods in MATLAB}, SIAM, Philadelphia, 2000.

\bibitem{LSY98}
R.~B. Lehoucq, D.~C. Sorensen, C.~Yang,
\emph{ARPACK Users' Guide}, SIAM, Philadelphia, 1998.

\bibitem{VK-NH}
\NoteAuthor,
The Vakhitov--Kolokolov criterion beyond Hermiticity,
submitted (2026); \VKarXiv.

\end{thebibliography}
\end{document}
